\section{PART I --- REAL ANALYSIS FOUNDATIONS}

\subsection{1. Metric Spaces and Topology}
\begin{itemize}
\item Metrics, norms, inner products; equivalence and comparison
\item Open and closed sets; interior, closure, boundary
\item Compactness: sequential, covering, Heine-Borel, and equivalence in metric spaces
\item Completeness; Baire category theorem and its consequences
\item Connectedness; path-connectedness
\end{itemize}

\subsection{2. Sequences and Function Spaces}
\begin{itemize}
\item Pointwise vs.\ uniform convergence; counterexamples
\item Equicontinuity; Arzel\`a-Ascoli theorem
\item Stone-Weierstrass theorem; polynomial approximation
\item Power series; radius of convergence; real analytic functions
\end{itemize}

\subsection{3. Multivariable Differentiation}
\begin{itemize}
\item Fr\'echet derivative; directional derivatives; their relationship
\item Chain rule; mean value inequality in $\mathbb{R}^n$
\item Inverse and implicit function theorems
\item Higher derivatives; Taylor's theorem; local extrema and second-order conditions
\end{itemize}

\subsection{4. Why a Better Integral is Needed}
\begin{itemize}
\item Riemann-Stieltjes integral; limitations
\item Failure modes of Riemann: pointwise limits, Fourier series
\item Informal preview of measure theory motivation
\item Key phenomena that force a better integral
\end{itemize}

\section{PART II --- MEASURE, INTEGRATION, AND PROBABILISTIC TOOLS}

\subsection{5. Measures and $\sigma$-Algebras}
\begin{itemize}
\item $\sigma$-algebras; measurable spaces; Borel sets
\item Measures; outer measures; Carath\'eodory extension theorem
\item Regularity properties; Borel and Radon measures on $\mathbb{R}^n$
\end{itemize}

\subsection{6. Lebesgue Integration}
\begin{itemize}
\item Measurable functions; simple functions and approximation
\item Lebesgue integral: definition through non-negative functions
\item Monotone convergence; Fatou's lemma; dominated convergence
\item Relation to Riemann; null sets and almost-everywhere statements
\end{itemize}

\subsection{7. $L^p$ Spaces}
\begin{itemize}
\item H\"older and Minkowski inequalities
\item Completeness: Riesz-Fischer theorem
\item Dense subsets: continuous and smooth approximation
\item Dual of $L^p$ $(1 < p < \infty)$; weak convergence in $L^p$
\end{itemize}

\subsection{8. Differentiation Theory}
\begin{itemize}
\item Hardy-Littlewood maximal function; weak $(1,1)$ estimate
\item Lebesgue differentiation theorem
\item Absolutely continuous functions; Radon-Nikodym theorem
\item Functions of bounded variation; Jordan decomposition
\end{itemize}

\subsection{9. Product and Signed Measures}
\begin{itemize}
\item Product $\sigma$-algebras and product measures
\item Fubini and Tonelli theorems; interchange of integrals
\item Hahn and Jordan decompositions
\item Radon-Nikodym theorem and Lebesgue-Radon-Nikodym decomposition
\item Conditional expectation as Radon-Nikodym derivative (preview)
\end{itemize}

\subsection{10. Concentration Inequalities}
\begin{itemize}
\item Markov and Chebyshev; moment generating functions
\item Sub-Gaussian random variables; Hoeffding's inequality
\item Bernstein's inequality; sub-exponential variables
\item Azuma-Hoeffding for martingales; McDiarmid's bounded differences method
\item Applications: generalization bounds in learning, MCMC mixing
\end{itemize}

\subsection{11. Large Deviations}
\begin{itemize}
\item Cram\'er's theorem on $\mathbb{R}$; rate function; Legendre-Fenchel connection
\item G\"artner-Ellis theorem; tilting
\item Sanov's theorem; relative entropy as rate function
\item Concentration as finite-sample control; large deviations as asymptotic control
\end{itemize}

\section{PART III --- FUNCTIONAL ANALYSIS AND FUNCTION SPACES}

\subsection{12. Normed Spaces and Banach Spaces}
\begin{itemize}
\item Norms; completeness; examples ($\ell^p$, $L^p$, $C(K)$, Sobolev)
\item Finite vs.\ infinite dimensions; Riesz lemma on non-compactness of unit ball
\item Baire category theorem in Banach spaces; generic properties
\item Series in Banach spaces; absolute vs.\ unconditional convergence
\end{itemize}

\subsection{13. Bounded Linear Operators}
\begin{itemize}
\item Operator norm; space $B(X,Y)$ and its completeness
\item Open mapping theorem and bounded inverse theorem
\item Closed graph theorem
\item Uniform boundedness principle (Banach-Steinhaus)
\end{itemize}

\subsection{14. Hilbert Spaces}
\begin{itemize}
\item Inner products; parallelogram law; polarization identity
\item Orthonormal bases; Gram-Schmidt; Parseval's identity
\item Projection theorem; orthogonal complements; best approximation
\item Riesz representation theorem; self-duality of $L^2$
\end{itemize}

\subsection{15. Dual Spaces and Weak Topologies}
\begin{itemize}
\item Hahn-Banach: analytic and geometric forms; separation
\item Dual spaces of $L^p$, $\ell^p$, $C(K)$
\item Reflexive spaces; bidual embedding
\item Weak and weak-* convergence; Banach-Alaoglu theorem
\end{itemize}

\subsection{16. Compact and Unbounded Operators}
\begin{itemize}
\item Compact operators; precompact sets; examples
\item Fredholm alternative; index theory (brief)
\item Spectral theorem for compact self-adjoint operators
\item Closed operators; adjoints; symmetric vs.\ self-adjoint
\item $C_0$-semigroups; Hille-Yosida theorem; evolution equations
\end{itemize}

\subsection{17. Distributions}
\begin{itemize}
\item Test functions $D(\Omega)$; mollifiers; regularization
\item Distributions as linear functionals; Dirac $\delta$ done properly
\item Differentiation and multiplication by smooth functions
\item Fundamental solutions; PDE examples (Laplacian, heat)
\end{itemize}

\subsection{18. Sobolev Spaces}
\begin{itemize}
\item Weak derivatives; $W^{k,p}$ and $H^k = W^{k,2}$ spaces
\item Sobolev embedding theorems (integer and fractional)
\item Trace theorem; extension operators
\item Poincar\'e inequality; Rellich-Kondrachov compactness
\end{itemize}

\section{PART IV --- FOURIER AND HARMONIC ANALYSIS}

\subsection{19. Fourier Analysis on $\mathbb{R}^n$}
\begin{itemize}
\item Fourier transform; Plancherel theorem; Hausdorff-Young
\item Schwartz space $S(\mathbb{R}^n)$ and tempered distributions
\item Fourier transform and regularity: $H^s$ Sobolev spaces
\item Paley-Wiener theorem; analytic signals
\end{itemize}

\subsection{20. Fourier Analysis on Groups}
\begin{itemize}
\item Locally compact abelian groups; Haar measure; Pontryagin duality
\item Fourier series on $\mathbb{T}^n$; discrete Fourier transform on $\mathbb{Z}/N\mathbb{Z}$
\item Peter-Weyl theorem for compact groups
\item Representation theory as generalized Fourier analysis
\end{itemize}

\subsection{21. Singular Integrals and Calder\'on-Zygmund Theory}
\begin{itemize}
\item Singular integrals: Hilbert transform, Riesz transforms
\item Calder\'on-Zygmund decomposition; good/bad parts
\item $L^p$ boundedness via weak-type $(1,1)$ and Marcinkiewicz interpolation
\item $T(1)$ theorem (statement); connections to elliptic PDE regularity
\end{itemize}

\subsection{22. Littlewood-Paley Theory and Wavelets}
\begin{itemize}
\item Dyadic decomposition; Littlewood-Paley square functions
\item Characterization of $L^p$ and $H^s$ norms via LP decomposition
\item Multiresolution analysis; Daubechies wavelets; vanishing moments
\item Fast wavelet transform; applications to signal processing and PDEs
\end{itemize}

\subsection{23. Hardy Spaces and BMO}
\begin{itemize}
\item Real Hardy space $H^1$ and its atomic decomposition
\item BMO: bounded mean oscillation; John-Nirenberg inequality
\item Duality $H^1 \leftrightarrow \mathrm{BMO}$
\item Endpoint estimates for singular integrals
\end{itemize}

\subsection{24. Interpolation and Spectral Methods}
\begin{itemize}
\item Riesz-Thorin complex interpolation
\item Marcinkiewicz weak-type interpolation
\item $K$-functional and real interpolation (brief)
\item Spectrum and resolvent; spectral radius formula
\item Spectral theorem and Borel functional calculus
\item Applications: wave propagation, Schr\"odinger equation
\end{itemize}

\section{PART V --- VARIATIONAL, CONVEX, AND TRANSPORT METHODS}

\subsection{25. Variational Methods and Calculus of Variations}
\begin{itemize}
\item Euler-Lagrange equations; derivation and examples
\item Direct methods: lower semicontinuity, coercivity, existence
\item Convex and polyconvex energies; relaxation
\item Regularity; connection to PDEs
\end{itemize}

\subsection{26. Convex Analysis}
\begin{itemize}
\item Convex sets and functions; epigraph; closure
\item Subdifferentials; Fenchel conjugate (Legendre-Fenchel transform)
\item Minimax theorems; saddle points; von Neumann
\item Variational inequalities; connection to optimization
\end{itemize}

\subsection{27. Optimal Transport}
\begin{itemize}
\item Monge problem; Kantorovich relaxation; linear programming formulation
\item Wasserstein distances; metrization of weak convergence
\item Brenier's theorem; polar factorization
\item Wasserstein gradient flows; JKO minimizing movement scheme
\item Connections to inverse problems, PDE, and generative modeling
\end{itemize}

\section{PART VI --- STOCHASTIC, HIGH-DIMENSIONAL, AND LEARNING-THEORETIC ANALYSIS}

\subsection{28. Stochastic Analysis}
\begin{itemize}
\item Brownian motion: construction, path properties, quadratic variation
\item Martingales in continuous time; optional stopping; Doob's inequalities
\item It\^o integral and It\^o's formula
\item SDEs: existence and uniqueness; Girsanov; Feynman-Kac formula
\end{itemize}

\subsection{29. Ergodic Theory}
\begin{itemize}
\item Measure-preserving transformations; ergodicity; invariant measures
\item Von Neumann $L^2$ ergodic theorem
\item Birkhoff pointwise ergodic theorem
\item Mixing conditions; spectral characterization
\item Applications: MCMC convergence; statistical mechanics; time series
\end{itemize}

\subsection{30. Random Matrix Theory}
\begin{itemize}
\item Wigner semicircle law: empirical spectral distribution of symmetric random matrices
\item Marchenko-Pastur law: spectrum of sample covariance matrices
\item Tracy-Widom distribution; universality (conceptual statement)
\item Applications: PCA in high dimensions; deep learning initialization; free probability
\end{itemize}

\subsection{31. Compressed Sensing and Sparse Recovery}
\begin{itemize}
\item Sparse signals; $\ell_1$-relaxation; basis pursuit
\item Restricted Isometry Property (RIP); incoherence
\item RIP from Gaussian and Fourier random matrices
\item Recovery guarantees; LASSO and statistical interpretation
\end{itemize}

\subsection{32. Empirical Processes and Learning Theory}
\begin{itemize}
\item Glivenko-Cantelli; uniform laws of large numbers
\item Rademacher complexity; connections to covering numbers
\item VC dimension and the VC inequality
\item PAC learning; generalization bounds from empirical process theory
\end{itemize}

\subsection{33. Neural Network Approximation Theory}
\begin{itemize}
\item Universal approximation: Cybenko, Hornik
\item Approximation rates for smooth functions; depth vs.\ width tradeoffs
\item Barron spaces; integral representation of shallow networks
\item Curse of dimensionality: what networks overcome and what they do not
\item Connections to ridge functions and harmonic analysis
\end{itemize}

\section{PART VII --- ANALYTIC TOOLBOX FOR APPLIED MATHEMATICS}

\subsection{34. Convolution, Kernels, and Smoothing}
\begin{itemize}
\item Convolution as localization and smoothing
\item Approximate identities and mollification
\item Heat kernel intuition and regularization mechanisms
\item Fourier-side interpretation of filtering
\end{itemize}

\subsection{35. Compactness and Weak Convergence Methods}
\begin{itemize}
\item Arzel\`a-Ascoli, Rellich, and Prokhorov as compactness tools
\item Weak convergence, weak-* convergence, and tightness
\item Passing to limits in nonlinear problems
\item Existence proofs in PDE, probability, and variational problems
\end{itemize}

\subsection{36. Energy, Duality, and Bootstrap Arguments}
\begin{itemize}
\item Energy estimates and a priori bounds
\item Duality arguments and transference of estimates
\item Bootstrap arguments for regularity improvement
\item Perturbation methods and fixed-point methods
\end{itemize}

\subsection{37. Reading Modern Applied Analysis}
\begin{itemize}
\item TT* arguments and interpolation tricks
\item Dyadic decompositions; soft vs.\ hard cutoffs
\item When to use Sobolev, Fourier, variational, or probabilistic tools
\item How modern papers combine analysis with numerics, optimization, and learning
\end{itemize}
